\paragraph{AQ-136-PLOCAL: common sources over every $\mathbb Z_{(p_0)}$-algebra.}
Fix a prime $p_0$.  If $k$ is any commutative
$\mathbb Z_{(p_0)}$-algebra, every finite family of admissible
strict PD theories in $\mathcal A_k$ has a common admissible
source.  In particular this covers $\mathbb Z_{p_0}$,
$\mathbb F_{p_0}$ and $\mathbb Z/p_0^r$ with their original
coefficients.  The assertion is common-source existence;
it does not here assert terminality of the $p_0$-local
all-arrow groupoid.

\begin{proof}
Put $S=\Gamma_k(T)$ and choose $N=p_0^e$, $e\geq0$.
The PD homomorphism
\[
 \rho_N:\Gamma_k(U)\longrightarrow\Gamma_k(T),
                  \qquad U\longmapsto T^{[N]},              \tag{136.21}
\]
becomes a homogeneous map if $U$ is assigned weight $N$.
For every $m\geq0$ it sends the basis element $U^{[m]}$ to
\[
 c_mT^{[mN]},\qquad c_m=\frac{(mN)!}{m!(N!)^m}.
                                                               \tag{136.22}
\]
The integer $c_m$ is a $p_0$-adic unit.  Indeed counting the
multiples of successive powers of $p_0$ in a factorial gives
\[
 v_{p_0}((mp_0^e)!)
  =m\frac{p_0^e-1}{p_0-1}+v_{p_0}(m!),
 \qquad
 v_{p_0}((p_0^e)!)=\frac{p_0^e-1}{p_0-1};
\]
the difference in (136.22) is zero.
For $0\leq a<N$, the integer
$\binom{mN+a}{a}$ is also a $p_0$-adic unit.
For powers $p_0^j$ with $j\leq e$, divisibility of $mN$
by $p_0^j$ makes the corresponding three factorial floors
cancel.  For $j>e$, $a<p_0^e$ implies
\[
 \left\lfloor\frac{mN+a}{p_0^j}\right\rfloor
  =\left\lfloor\frac{m}{p_0^{j-e}}\right\rfloor,
 \qquad \left\lfloor\frac a{p_0^j}\right\rfloor=0,
\]
so those floors also cancel.  These integers are therefore
units in every $\mathbb Z_{(p_0)}$-algebra $k$.

The products
\[
 T^{[a]}\rho_N(U^{[m]})
    =c_m\binom{mN+a}{a}T^{[mN+a]}
\]
are consequently unit multiples of all the distinct basis
elements of $S$.  This proves the actual graded-module
isomorphism
\[
 S=\bigoplus_{a=0}^{N-1}T^{[a]}\Gamma_k(U)                 \tag{136.23}
\]
via $\rho_N$.  In particular $S$ is free, not merely flat,
over the indicated divided-power coefficient algebra.
The claim holds even if $k$ has $p_0$-torsion or nilpotents.

Let $A$ be the admissible integral DP simplex theory
(136.5), with its coefficient named $U$, and regrade
all its weights by $N$.  Form
\[
 B_N^{\rm PD}=S\otimes_{\Gamma_k(U)}A .                    \tag{136.24}
\]
In each simplex this is exactly the PD quotient with
$\sum x_j=T^{[N]}$ and $\sum e_j=0$ from (136.18).
Eliminating $x_0,e_0$ identifies it with the free PD
polynomial de Rham complex over $S$ in the remaining
variables.  The integral contractions in (136.6) prove
vertical cohomology $S$ in degree zero and zero otherwise,
in every weight and every simplex.

For the horizontal calculation, (136.23) is a natural
direct sum decomposition of the whole simplicial module,
since its coefficients are constant simplicially.
The component of weight $w=Nm+a$, $0\leq a<N$, and
cochain degree $r$ is identified, as a simplicial
$k$-module, with $A_*^{m,r}$.  The latter is contractible
for $m>r$, by (136.7).  Thus
$\pi_{B_N^{\rm PD}}(w)\geq\lfloor w/N\rfloor\to\infty$,
and $B_N^{\rm PD}$ is admissible.  The coefficient tail
is complete: every $T^{[w]}$ occurs in the free basis
calculation, with no rationalization.

Given finitely many admissible targets $C_b$, choose
$e$ so large that every row $C_b^{N,0}$ is acyclic.
Condition (136.4) permits this along the unbounded
sequence $N=p_0^e$.  Lift the constant coefficient
$T^{[N]}$ along $k\to L$ to a simplicial linear map
$L\to C_b^{N,0}$, using exactly the normalized-disk
argument of (136.15).  Its coordinate images have
sum $T^{[N]}$.  Send $x_j$ to those images and
$e_j$ to their vertical differentials.  The PD
polynomial universal property, together with
$d\gamma_m(x)=\gamma_{m-1}(x)dx$, gives a strict
simplicial differential PD map
$B_N^{\rm PD}\to C_b$ over $S$.
It preserves both gradings and the two sum relations.
These are the required simultaneous source maps.
\end{proof}

\paragraph{Finite-product and positive-characteristic consequence.}
Common-source existence also holds for a finite product
$k=\prod_{i=1}^r k_i$ of rings covered by AQ-136-PLOCAL.
Indeed the coefficient idempotents split every
$k$-theory, its divided powers and its morphisms into
$r$ component theories over $k_i$.
Both vertical cohomology and normalized simplicial
homology respect this finite direct-product decomposition.
Thus admissibility is equivalent to admissibility of all
components, using the maximum of their finitely many
cutoffs.  Take componentwise common sources and their
finite product; this is the required admissible source
over $k$.  Zero component rings cause no exception.

In particular common-source existence holds for every
commutative ring of positive characteristic.  If its
characteristic is $m=\prod_i p_i^{a_i}>1$, the Chinese
remainder idempotents of $\mathbb Z/m$ give
$k=\prod_i k_i$, each a $\mathbb Z/p_i^{a_i}$-algebra
and hence a $\mathbb Z_{(p_i)}$-algebra.  The preceding
argument applies.  The zero ring is harmless as well.
This finite idempotent decomposition is not available
for $\mathbb Z$; no prime-varying integral coefficient
object is substituted for the single ring in that
remaining question.

\paragraph{AQ-136-PATH-OBSTRUCTION: the unweighted rational interval has no characteristic-$p$ replacement of the same kind.}
Let $k\ne0$ be a commutative ring of characteristic
$p_0>0$.  There is no nonnegative commutative differential
$k$-algebra $R$ with coefficient isomorphism
$k\simeq H^0(R)$, two $k$-algebra augmentations
$\epsilon_0,\epsilon_1:R\to k$, and an element
$s\in R^0$ such that $\epsilon_0(s)=0$ and
$\epsilon_1(s)=1$.

\begin{proof}
Since $R$ is nonnegative, $H^0(R)=\ker(d:R^0\to R^1)$
without a quotient by boundaries.  The Leibniz rule and
characteristic $p_0$ give
$d(s^{p_0})=p_0s^{p_0-1}ds=0$.
The coefficient isomorphism therefore forces
$s^{p_0}$ to be the image of some $c\in k$.
Applying the two augmentations gives both $c=0$ and
$c=1$, contradicting $k\ne0$.
\end{proof}

This last obstruction rules out only such an unweighted
strict interval.  Weighted comparisons and other
localization arguments are not excluded.  In particular
it is consistent with AQ-136-PLOCAL, and it supplies no
counterexample to the still separate integral or
$p_0$-local terminal-localization question.
